\section{\large Título I}

Si quiero citar dentro del texto según~\textcite{cervantes1999} 
o si al final del texto~\parencite{cervantes1999}. O una cita tradicional~\cite{cervantes1999} dentro del texto.

En un lugar de la Mancha, de cuyo nombre no quiero acordarme, no ha mucho tiempo que vivía un hidalgo de los de lanza en astillero, adarga antigua, rocín flaco y galgo corredor.
\subsection{Título II} 
Una olla de algo más vaca que carnero, salpicón las más noches, duelos y quebrantos los sábados, lantejas los viernes, algún palomino de añadidura los domingos, consumían las tres partes de su hacienda. 

\subsubsection{Título III}
El resto della concluían sayo de velarte, calzas de velludo para las fiestas, con sus pantuflos de lo mesmo, y los días de entresemana se honraba con su vellorí de lo más fino.
\paragraph{Título IV}
Tenía en su casa una ama que pasaba de los cuarenta, y una sobrina que no llegaba a los veinte, y un mozo de campo y plaza, que así ensillaba el rocín como tomaba la podadera.
\myparagraph{Título IV ii}
Frisaba la edad de nuestro hidalgo con los cincuenta años; era de complexión recia, seco de carnes, enjuto de rostro, gran madrugador y amigo de la caza. 
\subparagraph{Título V}
Quieren decir que tenía el sobrenombre de Quijada, o Quesada, que en esto hay alguna diferencia en los autores que deste caso escriben; aunque por conjeturas verosímiles se deja entender que se llamaba Quijana.

\lstinputlisting[
    style=python_style,
    linerange={11-},
    %firstnumber=11,
    caption={Calcular el área y perímetro de un círculo.},
    label={cod:punto1}
]{src/punto1.py}

\lipsum[1] %Agrega 3 párramos de lorem ipsum
\section{Cómo agregar imágenes}

En la figura~\ref{fig:placeholder} se observa una imagen que se puede referenciar dentro del texto.

\begin{figure}
    \caption{Caption}
    \label{fig:placeholder}
    \centering
    \includegraphics[width=0.9\linewidth]{imagen1}
\end{figure}




\lipsum[2-3] %Agrega 5 párramos de lorem ipsum

\section{Cómo agregar tablas}

En la tabla~\ref{tab:ejemplo_notas} se muestra cómo referenciar una tabla en \LaTeX{} dentro del texto.

\begin{table}[h]
\centering
\caption{Comparativa de rendimientos académicos}
\label{tab:ejemplo_notas}
\begin{tabular}{|l|c|r|}
\hline
\textbf{Estudiante} & \textbf{Asignatura} & \textbf{Nota Final} \\ \hline
Ana García          & Matemáticas         & 9.5                 \\ \hline
Luis Pérez          & Historia            & 8.2                 \\ \hline
Elena Rivas         & Física              & 10.0                \\ \hline
\end{tabular}
\end{table}

\lipsum[1-3]

\begin{figure}
	\caption{Varias Figuras}
	\begin{minipage}[c]{0.6\textwidth}
		\begin{minipage}[c]{0.48\textwidth}
			\centering
			\subcaption{Figura Uno}
			\includegraphics[width=0.95\textwidth]{imagen1}
			\label{fig:FiguraSuperiorIzquierda}
		\end{minipage}
		\begin{minipage}[c]{0.48\textwidth}
			\centering
			\subcaption{Figura Uno}
			\includegraphics[width=0.95\textwidth]{imagen1}
			\label{fig:fig:FiguraSuperiorCentro}
		\end{minipage}
		\begin{minipage}[c]{0.48\textwidth}
			\centering
			\subcaption{Figura Uno}
			\includegraphics[width=0.95\textwidth]{imagen1}
			\label{fig:FiguraInferiorIzquierda}
		\end{minipage}
		\begin{minipage}[c]{0.48\textwidth}
			\centering
			\subcaption{Figura Uno}
			\includegraphics[width=0.95\textwidth]{imagen1}
			\label{fig:FiguraInferiorInferiorCentro}
		\end{minipage}
	\end{minipage}
	\begin{minipage}[c]{0.32\textwidth}
		\centering
		\subcaption{Figura Cinco}
		\includegraphics[width=0.95\textwidth]{imagen2}
		\label{fig:FiguraDerecha}
	\end{minipage}
	\label{fig:Varias1}
\end{figure}

